\documentclass[UTF8]{ctexart}
\usepackage[a4paper,margin=2.5cm]{geometry}
\usepackage{amsmath,amssymb}
\usepackage{graphicx}
\usepackage{booktabs}
\usepackage{hyperref}

\title{第六题：MNIST 分类中 SGD 与自然梯度训练的对比}
\author{}
\date{\today}

\begin{document}
\emergencystretch=2em
\maketitle

\section{题目与实现说明}

第六题要求构建神经网络对 MNIST 数据集中的 $0$--$9$ 手写数字进行分类，并分别使用随机梯度下降
（SGD）和自然梯度（Natural Gradient, NG）训练，最后比较两种方法的训练结果。
本解答的程序为 \texttt{mnist\_sgd\_ng.py}。程序优先读取标准 MNIST IDX 文件，包括训练图像、
训练标签、测试图像和测试标签四个文件；文件名可以带 \texttt{.gz} 后缀。当前已经将标准 MNIST
IDX 文件下载到 \texttt{data/mnist}，因此本报告采用真实 MNIST 全量训练集和测试集进行实验。

为了进一步提高两种方法的最终准确率，本次在前一版完整 MNIST 实验基础上继续增加训练轮数：
训练集使用全部 $60000$ 个样本，测试集使用全部 $10000$ 个样本，隐藏层宽度设为 $128$，训练轮数
由 $20$ 轮增加到 $30$ 轮；同时采用较稳定的学习率组合，使后期准确率能够继续上升而不过早震荡。
运行命令为
\begin{quote}
\small
\begin{verbatim}
python mnist_sgd_ng.py --mnist-dir data/mnist --require-mnist
  --output-dir results_mnist_full_30ep
  --train-size 60000 --test-size 10000 --epochs 30
  --hidden-dim 128 --sgd-lr 0.08
  --ng-lr 0.03 --ng-damping 0.5 --max-update-norm 1.5
\end{verbatim}
\end{quote}

\section{神经网络模型}

采用一个单隐藏层多层感知机（MLP）：
\[
h=\mathrm{ReLU}(xW_1+b_1),\qquad
p_\theta(y\mid x)=\mathrm{softmax}(hW_2+b_2),
\]
其中输入 $x\in\mathbb{R}^{784}$，隐藏层维度为 $128$，输出层为 $10$ 类。损失函数为带 $L_2$
正则项的交叉熵：
\[
\mathcal{L}(\theta)=
-\frac{1}{n}\sum_{i=1}^{n}\log p_\theta(y_i\mid x_i)
\;+\;\frac{\lambda}{2}\left(\|W_1\|_2^2+\|W_2\|_2^2\right).
\]

\section{SGD 与自然梯度方法}

SGD 直接沿小批量梯度的反方向更新参数：
\[
\theta_{t+1}=\theta_t-\eta\,\hat g_t,
\]
其中 $\hat g_t=\nabla_\theta \mathcal{L}_B(\theta_t)$ 为当前 minibatch 的梯度。

自然梯度的思想是不使用欧氏距离衡量参数变化，而使用模型分布诱导出的 Fisher 信息矩阵
$F(\theta)$ 衡量参数空间的局部几何：
\[
\theta_{t+1}=\theta_t-\alpha\,F(\theta_t)^{-1}\hat g_t.
\]
对 MLP 精确构造和求逆完整 Fisher 矩阵代价很高，因此程序采用对角经验 Fisher 近似：
\[
\hat F_t=\beta \hat F_{t-1}+(1-\beta)(\hat g_t\odot \hat g_t),
\]
\[
\theta_{t+1}
=\theta_t-\alpha\left(\mathrm{diag}(\hat F_t)+\lambda_{\mathrm{NG}}I\right)^{-1}\hat g_t.
\]
其中 $\lambda_{\mathrm{NG}}$ 为阻尼项，用于避免 Fisher 对角元素过小导致步长过大。该实现保留了
自然梯度按参数敏感度自适应缩放梯度的核心思想，同时避免了完整 Fisher 矩阵的存储和求逆。

\section{实验设置}

两种优化器使用相同的数据划分、相同初始参数和相同 minibatch 顺序。主要参数如下：

\begin{center}
\begin{tabular}{ll}
\toprule
项目 & 数值 \\
\midrule
训练集大小 & 60000 \\
测试集大小 & 10000 \\
隐藏层神经元 & 128 \\
训练轮数 & 30 \\
Batch size & 128 \\
SGD 学习率 & 0.08 \\
NG 学习率 & 0.03 \\
NG 阻尼 & 0.5 \\
NG Fisher 衰减系数 & 0.9 \\
\bottomrule
\end{tabular}
\end{center}

\begin{figure}[htbp]
\centering

\includegraphics[width=0.72\linewidth]{results_mnist_full_30ep/sample_digits.png}
\caption{标准 MNIST 数据集中的 $28\times 28$ 手写数字样例}
\end{figure}

图 1 展示了真实 MNIST 实验中抽取的数字样例。每个样本是 $28\times 28$ 灰度图，数字存在明显的
书写风格、粗细、倾斜和局部形变差异，因此比离线模板数据更接近真实分类任务。

\section{运行结果对比}

本次 $30$ 轮训练的最终结果如下：

\begin{center}
\begin{tabular}{lccccc}
\toprule
优化器 & 最终训练准确率 & 最终测试准确率 & 最终测试损失 & 最佳测试准确率 & 最佳轮数 \\
\midrule
SGD & 0.9902 & 0.9771 & 0.0966 & 0.9774 & 29 \\
NG & 0.9863 & 0.9760 & 0.1028 & 0.9760 & 30 \\
\bottomrule
\end{tabular}
\end{center}

\begin{figure}[htbp]
\centering
\includegraphics[width=0.86\linewidth]{results_mnist_full_30ep/accuracy_curve.png}
\caption{SGD 与 NG 的测试准确率曲线}
\end{figure}

图 2 给出了测试集准确率随训练轮数变化的曲线。两种方法在前 $10$ 轮快速上升，之后进入缓慢提升阶段。
当训练轮数从 $20$ 轮增加到 $30$ 轮时，SGD 的最终测试准确率由 $0.9741$ 提高到 $0.9771$，
NG 的最终测试准确率由 $0.9717$ 提高到 $0.9760$。这说明前一版 $20$ 轮实验还没有完全达到平台期，
适当增加训练轮数确实能带来进一步提升。SGD 的最佳点出现在第 $29$ 轮，NG 的最佳点出现在第 $30$ 轮，
说明 NG 在该参数设置下后期仍保持小幅上升趋势。

\begin{figure}[htbp]
\centering
\begin{minipage}{0.48\linewidth}
\centering
\includegraphics[width=\linewidth]{results_mnist_full_30ep/confusion_sgd.png}
\caption{SGD 混淆矩阵}
\end{minipage}
\hfill
\begin{minipage}{0.48\linewidth}
\centering
\includegraphics[width=\linewidth]{results_mnist_full_30ep/confusion_ng.png}
\caption{NG 混淆矩阵}
\end{minipage}
\end{figure}

图 3 和图 4 为两种优化器在 $10000$ 个 MNIST 测试样本上的混淆矩阵。矩阵的行表示真实类别，列表示
预测类别；对角线越亮、非对角线数字越少，说明分类越准确。两种方法的大多数错误仍集中在形状相近的
类别之间，例如 $4$ 与 $9$、$3$ 与 $5$、$7$ 与 $9$ 等。增加训练轮数后，非对角线错误数量进一步减少，
说明模型对各类数字的决策边界更加稳定。SGD 的测试准确率略高，但 NG 的混淆矩阵整体结构与 SGD 接近，
表明对角 Fisher 近似同样学习到了主要类别边界。

\section{解释与讨论}

从实验结果可以看出，提高准确率最有效的操作包括三类。第一，使用完整训练集而不是小规模子集；
前期快速测试只使用 $4000$ 个训练样本和 $1000$ 个测试样本，SGD 与 NG 的测试准确率约为
$0.9070$ 和 $0.8820$，而全量 MNIST 实验已经将二者提升到 $97\%$ 以上。第二，适当增大隐藏层宽度；
隐藏层由 $64$ 增至 $128$ 后，网络具有更强的非线性表达能力。第三，在学习率稳定的前提下增加训练轮数；
从 $20$ 轮到 $30$ 轮，SGD 又提升约 $0.30$ 个百分点，NG 提升约 $0.43$ 个百分点。

两种方法的差异也值得说明。SGD 在当前单隐藏层 MLP 和固定学习率设置下略高于 NG，这是因为对角 Fisher
近似忽略了参数之间的相关性，不能完全等价于真实自然梯度。不过 NG 的后期曲线较平滑，且第 $30$ 轮仍是
最佳点，说明其自适应缩放在后期仍有继续优化的作用。若进一步引入学习率衰减、卷积神经网络、数据增强或
K-FAC 等更强的 Fisher 近似，准确率仍可能继续提升；但在本题的 NumPy 单隐藏层 MLP 框架下，当前结果已经
明显优于前一版快速实验，并能较公平地展示 SGD 与自然梯度近似的性能差异。

\section{文件说明}

\begin{itemize}
  \item \texttt{mnist\_sgd\_ng.py}：主程序，包含数据读取、MLP、SGD/NG 训练和结果输出；
  \item \texttt{requirements.txt}：运行依赖；
  \item \texttt{data/mnist/}：标准 MNIST IDX 数据文件；
  \item \texttt{results\_mnist\_full\_30ep/results.csv}：每轮训练与测试指标；
  \item \texttt{results\_mnist\_full\_30ep/summary.json}：最终对比摘要；
  \item \texttt{results\_mnist\_full\_30ep/accuracy\_curve.png}：准确率曲线；
  \item \texttt{results\_mnist\_full\_30ep/confusion\_sgd.png}：SGD 混淆矩阵；
  \item \texttt{results\_mnist\_full\_30ep/confusion\_ng.png}：NG 混淆矩阵。
\end{itemize}

\end{document}
